% Options for packages loaded elsewhere
\PassOptionsToPackage{unicode}{hyperref}
\PassOptionsToPackage{hyphens}{url}
\documentclass[
]{article}
\usepackage{xcolor}
\usepackage{amsmath,amssymb}
\setcounter{secnumdepth}{-\maxdimen} % remove section numbering
\usepackage{iftex}
\ifPDFTeX
  \usepackage[T1]{fontenc}
  \usepackage[utf8]{inputenc}
  \usepackage{textcomp} % provide euro and other symbols
\else % if luatex or xetex
  \usepackage{unicode-math} % this also loads fontspec
  \defaultfontfeatures{Scale=MatchLowercase}
  \defaultfontfeatures[\rmfamily]{Ligatures=TeX,Scale=1}
\fi
\usepackage{lmodern}
\ifPDFTeX\else
  % xetex/luatex font selection
\fi
% Use upquote if available, for straight quotes in verbatim environments
\IfFileExists{upquote.sty}{\usepackage{upquote}}{}
\IfFileExists{microtype.sty}{% use microtype if available
  \usepackage[]{microtype}
  \UseMicrotypeSet[protrusion]{basicmath} % disable protrusion for tt fonts
}{}
\makeatletter
\@ifundefined{KOMAClassName}{% if non-KOMA class
  \IfFileExists{parskip.sty}{%
    \usepackage{parskip}
  }{% else
    \setlength{\parindent}{0pt}
    \setlength{\parskip}{6pt plus 2pt minus 1pt}}
}{% if KOMA class
  \KOMAoptions{parskip=half}}
\makeatother
\setlength{\emergencystretch}{3em} % prevent overfull lines
\providecommand{\tightlist}{%
  \setlength{\itemsep}{0pt}\setlength{\parskip}{0pt}}
\usepackage{bookmark}
\IfFileExists{xurl.sty}{\usepackage{xurl}}{} % add URL line breaks if available
\urlstyle{same}
\hypersetup{
  hidelinks,
  pdfcreator={LaTeX via pandoc}}

\author{}
\date{}

\begin{document}

\section{Implicit Regularization in Deep Linear
Networks}\label{implicit-regularization-in-deep-linear-networks}

\subsection{Exercise Session --- 1h30, pen and
paper}\label{exercise-session-1h30-pen-and-paper}
Use LLMs to help you out (but make sure that you understand).
\textbf{References}:

\begin{itemize}
\tightlist
\item
  Saxe, McClelland \& Ganguli, \emph{Exact solutions to the nonlinear
  dynamics of learning in deep linear neural networks}, ICLR 2014 / PNAS
  2019
\item
  Achour, Malgouyres \& Gerchinovitz, \emph{The loss landscape of deep
  linear neural networks: a second-order analysis}, JMLR 2024
\item
  Tu, Aranguri \& Jacot, \emph{Mixed Dynamics In Linear Networks:
  Unifying the Lazy and Active Regimes}, 2024
\end{itemize}

\begin{center}\rule{0.5\linewidth}{0.5pt}\end{center}

\subsection{Setup and notation}\label{setup-and-notation}

We study deep linear networks (DLNs) with \(L\) weight matrices
\(W_1 \in \mathbb{R}^{d_1 \times d_0}, W_2 \in \mathbb{R}^{d_2 \times d_1}, \ldots, W_L \in \mathbb{R}^{d_L \times d_{L-1}}\).
The network computes:

\[f(x) = W_L W_{L-1} \cdots W_1 \, x =: W x\]

where \(W = W_L \cdots W_1 \in \mathbb{R}^{d_L \times d_0}\) is the
end-to-end (or ``student'') matrix. We train on a dataset
\(\{(x_\mu, y_\mu)\}_{\mu=1}^N\) with the squared loss:

\[\mathcal{L}(\theta) = \frac{1}{2N} \sum_{\mu=1}^{N} \|y_\mu - W_L \cdots W_1 \, x_\mu\|^2\]

In the population limit with whitened inputs \(\Sigma_X = I\), this
becomes:

\[\mathcal{L}(W) = \frac{1}{2}\|M - W\|_F^2\]

where \(M = \Sigma_{YX} \Sigma_X^{-1}\) is the ``teacher'' matrix (the
OLS solution) with SVD
\(M = U \,\text{diag}(s_1, \ldots, s_r, 0, \ldots, 0)\, V^\top\), and
\(s_1 \geq s_2 \geq \cdots \geq s_r > 0\).

\begin{center}\rule{0.5\linewidth}{0.5pt}\end{center}

\subsection{Problem 1 --- Loss Landscape Geometry (15
min)}\label{problem-1-loss-landscape-geometry-15-min}

This problem explores the critical point structure of deep linear
networks, following Achour et al.~(2024).

\textbf{(a)} Consider the \textbf{diagonal} case:
\(d_0 = d_1 = \cdots = d_L = d\), and the teacher is diagonal,
\(M = \text{diag}(s_1, \ldots, s_d)\), with
\(s_1 > s_2 > \cdots > s_d > 0\). Restrict attention to diagonal weight
matrices \(W_l = \text{diag}(w_l^{(1)}, \ldots, w_l^{(d)})\). Show that
the loss decomposes into \(d\) independent scalar problems:

\[\mathcal{L} = \frac{1}{2}\sum_{\alpha=1}^d \left(s_\alpha - \prod_{l=1}^L w_l^{(\alpha)}\right)^2\]

\textbf{(b)} For a single scalar mode with target \(s > 0\), find all
first-order critical points of
\(\ell(w_1, w_2) = \frac{1}{2}(s - w_1 w_2)^2\) at depth \(L = 2\). Show
that the critical points are:

\begin{itemize}
\tightlist
\item
  The global minimum manifold: \(w_1 w_2 = s\)
\item
  The origin: \(w_1 = w_2 = 0\)
\end{itemize}

Classify the origin as a saddle point by computing the Hessian \(H\) of
\(\ell\) at \((0,0)\) and showing it has both positive and negative
eigenvalues.

\emph{Hint}: Write the gradient equations 

\textbf{(c)} Achour et al.~show that every
first-order critical point \(\theta = (W_1, \ldots, W_L)\) satisfies:
there exists a subset \(S \subseteq \{1, \ldots, d\}\) such that

\[W = W_L \cdots W_1 = P_S\, M\]

where \(P_S = U_S U_S^\top\) is the orthogonal projector onto the span
of the left singular vectors \(\{u_\alpha\}_{\alpha \in S}\) of \(M\).
For the diagonal case \(M = \text{diag}(s_1, \ldots, s_d)\), write the
value of the loss at the critical point corresponding to
\(S = \{1, \ldots, r\}\) with \(r < d\).

\textbf{(d)} The group
\(GL_h := GL_{d_1} \times \cdots \times GL_{d_{L-1}}\) acts on the
weights by:

\[(W_1, \ldots, W_L) \mapsto (g_1 W_1,\; g_2 W_2 g_1^{-1},\; \ldots,\; W_L g_{L-1}^{-1})\]

Verify that the student map \(\mu(\theta) = W_L \cdots W_1\) is
invariant under this action: \(\mu(g \cdot \theta) = \mu(\theta)\). What
does this imply about the dimension of the set of global minima in
parameter space?

\begin{center}\rule{0.5\linewidth}{0.5pt}\end{center}

\subsection{Problem 2 --- Gradient Flow and Conserved Quantities (20
min)}\label{problem-2-gradient-flow-and-conserved-quantities-20-min}

We now study gradient flow:
\(\dot{\theta}(t) = -\nabla_\theta \mathcal{L}(\theta(t))\), the
continuous-time limit of gradient descent with infinitesimal learning
rate.

\textbf{(a)} For a two-layer diagonal DLN (\(L=2\)) with loss
\(\mathcal{L} = \frac{1}{2}\|M - W_2 W_1\|_F^2\), assume that each weight matrix $W_i$ is a diagonal matrix and whitened inputs.
Derive the gradient flow equations for each layer. Show that:

\[\dot{W}_1 = W_2^\top(M - W_2 W_1), \qquad \dot{W}_2 = (M - W_2 W_1) W_1^\top\]

\textbf{(b)} Define the \textbf{balancedness matrix}:

\[G := W_2^\top W_2 - W_1 W_1^\top\]

Show that \(G\) is conserved under the gradient flow,
i.e.~\(\dot{G} = 0\). In other words, gradient flow is constrained on the balanced manifold \begin{equation*}
    \mathcal{M}_{\beta}
    :=
    \{\theta\in\Omega\mid
    W_{\ell+1}^{\top}W_{\ell+1}
    -
    W_{\ell}W_{\ell}^{\top}
    =
    \beta_\ell,\;
    \ell=1,\ldots,L-1\}.
\end{equation*}

\emph{Hint}: Compute
\(\dot{G}\),
substitute the gradient flow equations and verify that the
terms cancel pairwise.

\textbf{(c)} From now on, assume \textbf{balanced initialization}:
\(G(0) = 0\), which by part (b) means \(W_2^\top W_2 = W_1 W_1^\top\)
for all time. We want to derive the gradient flow in function space (the
ODE for the student \(W = W_2 W_1\)). This requires several steps.

\textbf{(c.i)} Compute \(\dot{W} := \frac{d}{dt}(W_2 W_1)\) using the
gradient flow equations from (a). Show that:

\[\dot{W} = (M - W)\, W_1^\top W_1 + W_2 W_2^\top\, (M - W)\]

\textbf{(c.ii)} We now need to express \(W_1^\top W_1\) and
\(W_2 W_2^\top\) in terms of \(W = W_2 W_1\). For simplicity, assume
that the weight matrices \(W_l\) are diagonal. Show that:

\[W_1^\top W_1 = (W^\top W)^{1/2}\]

\textbf{(c.iii)} Similarly, show that
\(W_2 W_2^\top = (W W^\top)^{1/2}\).

\textbf{(c.iv)} Substitute the results of (c.ii) and (c.iii) into (c.i)
to obtain:

\[\dot{W} = (W W^\top)^{1/2}(M - W) + (M - W)(W^\top W)^{1/2}\]

\textbf{(d)} Define the NTK operator for \(L = 2\) as:

\[K[F] := (WW^\top)^{1/2}\, F + F\, (W^\top W)^{1/2}\]

Show that the gradient flow from (c) can be written as
\(\dot{W} = K[M - W]\).

It turns out that this result generalizes to depth \(L\) on the balanced
manifold and to non-diagonal weight matrices:

\[\dot{W} = \sum_{k=1}^{L} (WW^\top)^{\frac{L-k}{L}} (M - W) (W^\top W)^{\frac{k-1}{L}}\]
This is the NTK equation in the case of DLNs. The NTK equation is a gradient flow in function space with NTK being a preconditioning operator for the gradient.
\begin{center}\rule{0.5\linewidth}{0.5pt}\end{center}

\subsection{Problem 3 --- Rich Regime: saddle-to-saddle (25
min)}\label{problem-3-rich-regime-incremental-learning-25-min}

This is the core problem. We derive the exact solution of the rich
regime dynamics directly from the self-consistent equation of Problem 2,
emphasizing the NTK perspective.

\textbf{(a)} \textbf{Alignment assumption.} Start from the balanced
gradient flow equation derived in Problem 2(c):

\[\dot{W} = (WW^\top)^{1/2}(M - W) + (M - W)(W^\top W)^{1/2}\]

We work in the \textbf{rich regime} (small initialization) with balanced
weights. Assume that the student is \textbf{aligned} to the teacher: the
singular vectors of \(W(t)\) coincide with those of \(M\) at all times.
Concretely, write:

\[M = U\,\text{diag}(s_1, \ldots, s_d)\,V^\top, \qquad W(t) = U\,\text{diag}(w_1(t), \ldots, w_d(t))\,V^\top\]

where \(U, V\) are the left/right singular vectors of \(M\), and
\(w_\alpha(t) \geq 0\) are the evolving singular values of \(W\).

\textbf{(a.i)} Under this alignment assumption, show that the NTK
operator is aligned to the task, i.e.~\(K[F]\) preserves the SVD basis.
In particular, show that:

\[(WW^\top)^{1/2} = U\,\text{diag}(w_1, \ldots, w_d)\,U^\top, \qquad (W^\top W)^{1/2} = V\,\text{diag}(w_1, \ldots, w_d)\,V^\top\]

and that the residual is
\(M - W = U\,\text{diag}(s_1 - w_1, \ldots, s_d - w_d)\,V^\top\).

\emph{Hint}: Recall that for \(A = U\,\text{diag}(\sigma_i)\,V^\top\),
we have \(AA^\top = U\,\text{diag}(\sigma_i^2)\,U^\top\), and therefore
\((AA^\top)^{1/2} = U\,\text{diag}(|\sigma_i|)\,U^\top\).

\textbf{(a.ii)} Substitute into the self-consistent equation. Show that
the matrix equation decouples into \(d\) independent scalar ODEs:

\[\dot{w}_\alpha = 2\,w_\alpha\,(s_\alpha - w_\alpha), \qquad \alpha = 1, \ldots, d\]

\emph{Hint}: Compute the product \((WW^\top)^{1/2}(M - W)\) in the SVD
basis. It is diagonal with entries \(w_\alpha(s_\alpha - w_\alpha)\).
The second term contributes identically, giving the factor of 2.

\textbf{(b)} \textbf{Timescale of learning.} The ODE
\(\dot{w} = 2w(s - w)\) is a logistic equation whose solution is:

\[w(t) = \frac{s}{1 + \left(\frac{s}{w_0} - 1\right)e^{-2st}}\]

where \(w_0 = w(0)\). Compute the time \(t_\alpha\) it takes for mode
\(\alpha\) to travel from initial strength \(w_0\) to a final strength
\(w_f\). Show that:

\[t_\alpha = \frac{1}{2 s_\alpha} \ln\left(\frac{w_f\,(s_\alpha - w_0)}{w_0\,(s_\alpha - w_f)}\right)\]

Deduce that modes with larger singular values \(s_\alpha\) are learned
faster. This is a \textbf{separation of timescales}: the network learns
features in decreasing order of their singular value strength.

\textbf{(c)} \textbf{Incremental learning.} Consider a teacher with
\(r\) nonzero singular values and a small uniform initialization
\(w_\alpha(0) = w_0 \ll s_r\) for all \(\alpha\). Discuss:

\begin{itemize}
\tightlist
\item
  What is the approximate rank of \(W(t)\) at early times?
\item
  How does the rank evolve at intermediate times?
\item
  What happens in the limit \(t \to \infty\)?
\end{itemize}

What does this tell us about the sequence of critical points visited by
the gradient flow?

\begin{center}\rule{0.5\linewidth}{0.5pt}\end{center}

\subsection{Problem 4 --- Lazy Regime (15
min)}\label{problem-4-lazy-regime-15-min}

We now show that large initialization freezes the NTK and eliminates the
timescale separation found in Problem 3. For clarity, we work with the
\textbf{diagonal scalar} model from Problem 1(a), so each mode
\(\alpha\) is independent. This avoids matrix algebra and isolates the
essential mechanism.

\textbf{Setup.} Consider a single mode: a depth-2 diagonal
DLN with scalar weights \(a, b\) learning a target \(s > 0\). From
Problems 2--3, the balanced gradient flow for \(w = ab\) is:

\[\dot{w} = 2w(s - w)\]

The factor \(2w\) is the (scalar) NTK --- it is
\textbf{state-dependent}: the effective learning rate depends on the
current value of \(w\).

Now initialize at a \textbf{large} value \(w_0 \gg s > 0\). Assume that in
the early phase of training, while \(w\) has not changed much from
\(w_0\), the ODE becomes approximately:

\[\dot{w} \approx 2w_0(s - w)\]

\textbf{(b)} \textbf{Frozen-NTK solution.} Solve the linearized ODE from
(a). Show that:

\[w(t) = s + (w_0 - s)\,e^{-2w_0 t}\]

What is the learning timescale?

\textbf{(c)} \textbf{No timescale separation.} Now restore the mode
index \(\alpha\). In the lazy regime, each mode satisfies
\(\dot{w}_\alpha \approx 2w_0(s_\alpha - w_\alpha)\) with the
\textbf{same} rate \(2w_0\) for all \(\alpha\). Compare the lazy
learning time \(t_\alpha^{\text{lazy}}\) with the rich-regime timescale
from Problem 3(b): \(t_\alpha^{\text{rich}} \propto 1/s_\alpha\).

Discuss the difference in \textbf{rank bias} between the lazy and rich
regimes.

\begin{center}\rule{0.5\linewidth}{0.5pt}\end{center}

\subsection{Problem 5 --- Mixed Dynamics: Unifying Lazy and Rich (15
min)}\label{problem-5-mixed-dynamics-unifying-lazy-and-rich-15-min}

In practice, networks are neither purely lazy nor purely rich. Tu,
Aranguri \& Jacot (2024) provide a unified description. We develop the
key ideas using the diagonal scalar model.

\textbf{(a)} \textbf{The interpolating ODE.} In Problems 3 and 4, we
studied the same ODE \(\dot{w}_\alpha = 2w_\alpha(s_\alpha - w_\alpha)\)
in two limits: small \(w_0\) (rich) and large \(w_0\) (lazy). A more
general model for a depth-2 DLN with balanced initialization at scale
\(\sigma\) and width \(n\) replaces the scalar NTK \(2w_\alpha\) by:

\[\dot{w}_\alpha = 2\sqrt{w_\alpha^2 + \tau^2}\;(s_\alpha - w_\alpha)\]

where \(\tau := \sigma^2 \sqrt{n}\) is a \textbf{threshold} parameter
that depends on the initialization scale and width.

Verify that this ODE reproduces the two known regimes:

\begin{itemize}
\tightlist
\item
  \textbf{Rich limit} (\(\tau \to 0\)): recover
  \(\dot{w}_\alpha = 2|w_\alpha|(s_\alpha - w_\alpha)\).
\item
  \textbf{Lazy limit} (\(\tau \to \infty\) with \(w_\alpha\) bounded):
  recover \(\dot{w}_\alpha \approx 2\tau(s_\alpha - w_\alpha)\).
\end{itemize}

\textbf{(b)} \textbf{Two-phase dynamics.} At initialization, all modes
start at \(w_\alpha(0) \sim \sigma^2 \ll \tau\) (since
\(\sqrt{n} \gg 1\)). Observe that:

\begin{itemize}
\tightlist
\item
  When \(|w_\alpha| \ll \tau\), the effective learning rate is
  \(\approx 2\tau\), the same for all modes (lazy behavior).
\item
  When \(|w_\alpha| \gg \tau\), the effective learning rate is
  \(\approx 2|w_\alpha|\), which is mode-dependent (rich behavior).
\end{itemize}

Describe the resulting two-phase dynamics in words: what happens first,
and what happens later?

\textbf{(c)} \textbf{Mode-by-mode transition.} Not all modes cross the
threshold \(\tau\) at the same time: which modes cross it first? Explain
why this means that a single network can simultaneously have some modes
in the rich regime and others still in the lazy regime.

\begin{center}\rule{0.5\linewidth}{0.5pt}\end{center}

\subsection{Problem 6 (Bonus) --- Stochastic Implicit Bias (10
min)}\label{problem-6-bonus-stochastic-implicit-bias-10-min}

SGD introduces noise from mini-batching. Write the SGD update as a combination of gradient and noise term (use a LLM to check your answer). A continuous model is the
Langevin SDE:

\[d\theta_t = -\nabla \mathcal{L}(\theta_t)\,dt + \sqrt{\eta\,\Sigma(\theta_t)}\,dB_t\]

where \(\Sigma(\theta)\) is the covariance of the stochastic gradient
noise and \(\eta\) is the learning rate.

\textbf{(a)} The probability density \(p(\theta, t)\) of the parameters
evolves according to the Fokker--Planck equation:

\[\partial_t p = -\nabla \cdot \mathbf{j}, \qquad \mathbf{j} = -\nabla \mathcal{L}(\theta)\,p(\theta) - \frac{\eta}{2}\nabla \cdot \left(\Sigma(\theta)\,p(\theta)\right)\]

where \(\mathbf{j}\) is the probability current. Assume: (i)
stationarity \(\partial_t p^* = 0\), (ii) thermal equilibrium
\(\mathbf{j} = 0\), and (iii) isotropic noise \(\Sigma = \sigma^2 I\).
Show that the equilibrium distribution is the Boltzmann distribution:

\[p^*(\theta) \propto \exp\left(-\frac{2}{\eta \sigma^2}\mathcal{L}(\theta)\right)\]

\emph{Hint}: Setting \(\mathbf{j} = 0\) with \(\Sigma = \sigma^2 I\)
gives \(\nabla \mathcal{L}\, p + \frac{\eta\sigma^2}{2}\nabla p = 0\).
This is a first-order ODE for \(p\) in terms of \(\mathcal{L}\). Try the
ansatz \(p \propto e^{-\beta \mathcal{L}}\) and solve for \(\beta\).

\textbf{(b)} The ratio \(\frac{2}{\eta \sigma^2}\) plays the role of an
inverse temperature \(\beta\). Interpret what happens to the equilibrium
distribution when:

\begin{itemize}
\tightlist
\item
  \(\eta\) is very small (low temperature)
\item
  \(\eta\) is very large (high temperature)
\end{itemize}

Which regime favors flatter minima, and why might this be beneficial for
generalization?

\textbf{(c)} In practice, SGD noise is \textbf{not} isotropic:
\(\Sigma(\theta)\) depends on both the loss landscape and the data.
Without solving anything, explain qualitatively why anisotropic noise
can introduce an implicit bias that goes beyond what the loss function
\(\mathcal{L}\) alone would select. Specifically, why might SGD
preferentially escape sharp directions of the loss while remaining
stable along flat directions?

\end{document}
